\chapter{First results on rigid local systems}
\label{chap:1}
\setcounter{section}{-1}

\section{Generalities concerning rigid local systems over \texorpdfstring{\(\mathbb{C}\)}{C}}
\label{sec:1.0}

\sourcepara{1.0.1}
Let \(X\) be a projective smooth connected curve over \(\mathbb{C}\), of genus
\(g\), \(S\) a nonempty finite subset of \(X(\mathbb{C})\), and
\(U:=X-S\) the open complement. Suppose we are given on the complex manifold
\(U^{\mathrm{an}}\) a local system \(\mathcal{F}\), i.e., a locally constant
sheaf of finite-dimensional \(\mathbb{C}\)-vector spaces. As on any connected
complex manifold, if we fix a base point \(u\) in \(U^{\mathrm{an}}\), the
functor ``fibre at \(u\)'', \(\mathcal{F}\mapsto\mathcal{F}_u\), defines an
equivalence of categories
\[
  \text{local systems on }U^{\mathrm{an}}
  \cong
  \text{fin.-dim'l. }\mathbb{C}\text{-rep.'s of }\pi_1(U^{\mathrm{an}},u).
\]
We say that the local system \(\mathcal{F}\) is irreducible if the
corresponding representation \(\Lambda_{\mathcal{F}}\) of
\(\pi_1(U^{\mathrm{an}},u)\) is irreducible.

\sourcepara{1.0.2}
For every ``point at \(\infty\)'' \(s\) in \(S:=X-U\), the punctured
neighborhood
\[
  D^*(s):=U^{\mathrm{an}}\cap
  (\text{a small disc around }s\text{ in }X^{\mathrm{an}})
\]
is a punctured disc, whose fundamental group
\[
  I(s):=\pi_1(D^*(s),\text{ any base point})
\]
is canonically \(\mathbb{Z}\), with generator
\(\gamma_s:=\) ``turning once around \(s\) in the counterclockwise direction'',
the ``local monodromy transformation at \(s\)''. We say that two local systems
\(\mathcal{F}\) and \(\mathcal{G}\) on \(U^{\mathrm{an}}\) have isomorphic
local monodromy if for every ``point at \(\infty\)'' \(s\) in \(S:=X-U\),
there exists an isomorphism of local systems on \(D^*(s)\)
\[
  \mathcal{F}\vert D^*(s)\cong \mathcal{G}\vert D^*(s).
\]

\sourcepara{1.0.3}
We say that a local system \(\mathcal{F}\) on \(U^{\mathrm{an}}\) is
physically rigid if for every local system \(\mathcal{G}\) on
\(U^{\mathrm{an}}\) such that \(\mathcal{F}\) and \(\mathcal{G}\) have
isomorphic local monodromy, there exists an isomorphism
\(\mathcal{F}\cong\mathcal{G}\) of local systems on \(U^{\mathrm{an}}\).
Because we have assumed that \(S\) is nonempty, if \(\mathcal{F}\) and
\(\mathcal{G}\) have isomorphic local monodromy, they necessarily have the same
rank.

\sourcepara{1.0.4}
This notion of physical rigidity is reasonable only for genus zero. Indeed, if
\(X\) has genus \(g\geq 1\), there exist local systems \(\mathcal{L}\) of rank
one on \(X^{\mathrm{an}}\) no tensor power of which is trivial. (A rank one
local system on \(X^{\mathrm{an}}\) is a homomorphism from
\(\pi_1(X^{\mathrm{an}})^{\mathrm{ab}}\cong \mathbb{Z}^{2g}\) to
\(\mathbb{C}^{\times}\).) Denote by \(j:U^{\mathrm{an}}\to X^{\mathrm{an}}\)
the inclusion. Because the map
\[
  j_*:\pi_1(U^{\mathrm{an}},u)\to\pi_1(X^{\mathrm{an}},u)
\]
is surjective, no tensor power of \(j^*\mathcal{L}\) is trivial. But
\(j^*\mathcal{L}\) has trivial local monodromy, so for any local system
\(\mathcal{F}\) on \(U^{\mathrm{an}}\), \(\mathcal{F}\) and
\(\mathcal{F}\otimes j^*\mathcal{L}\) have isomorphic local monodromy. But
\(\mathcal{F}\) and \(\mathcal{F}\otimes j^*\mathcal{L}\) are not isomorphic
unless \(\mathcal{F}=0\), since already their determinants,
\(\det(\mathcal{F})\) and
\(\det(\mathcal{F})\otimes(j^*\mathcal{L})^{\otimes\operatorname{rank}(\mathcal{F})}\),
have a ratio which is nontrivial, indeed of infinite order. Thus no non-zero
local system \(\mathcal{F}\) on \(U^{\mathrm{an}}\) is physically rigid when
\(X\) has genus \(g\geq 1\).

\section{The case of genus zero}
\label{sec:1.1}

\sourcepara{1.1.1}
Let us now explore in greater detail the situation when \(X\) is
\(\mathbb{P}^1\). Even here, the situation is only understood for local systems
\(\mathcal{F}\) on \(U^{\mathrm{an}}\) which are irreducible. In this case,
there is a numerical criterion for physical rigidity.

\setcounter{theorem}{1}
\begin{theorem}
\label{thm:1.1.2}
Let \(S\) be a nonempty finite subset of \(\mathbb{P}^1(\mathbb{C})\), and
\(U:=\mathbb{P}^1-S\) the open complement,
\(j:U^{\mathrm{an}}\to(\mathbb{P}^1)^{\mathrm{an}}\) the inclusion,
\(\mathcal{F}\) an irreducible local system on \(U^{\mathrm{an}}\) of rank
\(n\geq 1\). Then \(\mathcal{F}\) is physically rigid if and only if
\[
  \chi\bigl((\mathbb{P}^1)^{\mathrm{an}},
  j_*\operatorname{End}(\mathcal{F})\bigr)=2.
\]
\end{theorem}

\begin{proof}
Suppose first that
\(\chi((\mathbb{P}^1)^{\mathrm{an}},j_*\operatorname{End}(\mathcal{F}))=2\),
and let \(\mathcal{G}\) be a local system on \(U^{\mathrm{an}}\) such that
\(\mathcal{F}\) and \(\mathcal{G}\) have isomorphic local monodromy. We will
show the existence of an isomorphism from \(\mathcal{F}\) to \(\mathcal{G}\).

For any local system \(\mathcal{H}\) on \(U^{\mathrm{an}}\), the
Euler--Poincaré formula states
\[
  \chi\bigl((\mathbb{P}^1)^{\mathrm{an}},j_*\mathcal{H}\bigr)
  =
  \chi(U^{\mathrm{an}},\mathbb{C})\operatorname{rank}(\mathcal{H})
  +\sum_{s\in S}\dim_{\mathbb{C}}\mathcal{H}^{I(s)}.
\]
Therefore if \(\mathcal{H}_1\) and \(\mathcal{H}_2\) are two local systems on
\(U^{\mathrm{an}}\) with isomorphic local monodromy, we have
\[
  \chi\bigl((\mathbb{P}^1)^{\mathrm{an}},j_*\mathcal{H}_1\bigr)
  =
  \chi\bigl((\mathbb{P}^1)^{\mathrm{an}},j_*\mathcal{H}_2\bigr).
\]
We apply this to
\(\mathcal{H}_1=\operatorname{End}(\mathcal{F})\) and to
\(\mathcal{H}_2=\operatorname{Hom}(\mathcal{F},\mathcal{G})\), which have
isomorphic local monodromy. Thus we find
\[
  \chi\bigl((\mathbb{P}^1)^{\mathrm{an}},
    j_*\operatorname{Hom}(\mathcal{F},\mathcal{G})\bigr)
  =
  \chi\bigl((\mathbb{P}^1)^{\mathrm{an}},
    j_*\operatorname{End}(\mathcal{F})\bigr)=2.
\]
But on a curve, \(\chi\) is \(h^0-h^1+h^2\leq h^0+h^2\), so we find
\[
  h^0\bigl((\mathbb{P}^1)^{\mathrm{an}},
    j_*\operatorname{Hom}(\mathcal{F},\mathcal{G})\bigr)
  +
  h^2\bigl((\mathbb{P}^1)^{\mathrm{an}},
    j_*\operatorname{Hom}(\mathcal{F},\mathcal{G})\bigr)
  \geq 2.
\]
We rewrite this in terms of ordinary and compact cohomology on
\(U^{\mathrm{an}}\) as
\[
  h^0\bigl(U^{\mathrm{an}},\operatorname{Hom}(\mathcal{F},\mathcal{G})\bigr)
  +
  h^2_c\bigl(U^{\mathrm{an}},\operatorname{Hom}(\mathcal{F},\mathcal{G})\bigr)
  \geq 2.
\]
The Poincaré dual of
\(\mathrm{H}^2_c(U^{\mathrm{an}},\operatorname{Hom}(\mathcal{F},\mathcal{G}))\)
is
\(\mathrm{H}^0(U^{\mathrm{an}},\operatorname{Hom}(\mathcal{G},\mathcal{F}))\),
so we obtain
\[
  h^0\bigl(U^{\mathrm{an}},\operatorname{Hom}(\mathcal{F},\mathcal{G})\bigr)
  +
  h^0\bigl(U^{\mathrm{an}},\operatorname{Hom}(\mathcal{G},\mathcal{F})\bigr)
  \geq 2.
\]
So at least one of the two groups
\[
  \mathrm{H}^0(U^{\mathrm{an}},\operatorname{Hom}(\mathcal{F},\mathcal{G}))
  =
  \operatorname{Hom}(\mathcal{F},\mathcal{G})
\]
or
\[
  \mathrm{H}^0(U^{\mathrm{an}},\operatorname{Hom}(\mathcal{G},\mathcal{F}))
  =
  \operatorname{Hom}(\mathcal{G},\mathcal{F})
\]
is nonzero. Since \(\mathcal{F}\) is irreducible and both
\(\mathcal{F}\) and \(\mathcal{G}\) have the same rank, any nonzero element of
either \(\operatorname{Hom}(\mathcal{F},\mathcal{G})\) or
\(\operatorname{Hom}(\mathcal{G},\mathcal{F})\) is necessarily an isomorphism.

Now suppose that \(\mathcal{F}\) is an irreducible local system of rank
\(n\geq 1\), which is physically rigid. We will show that
\[
  \chi\bigl((\mathbb{P}^1)^{\mathrm{an}},
    j_*\operatorname{End}(\mathcal{F})\bigr)=2.
\]
It suffices to show that for any local system \(\mathcal{F}\) of rank
\(n\geq 1\) which is physically rigid, we have
\[
  \chi\bigl((\mathbb{P}^1)^{\mathrm{an}},
    j_*\operatorname{End}(\mathcal{F})\bigr)\geq 2.
\]
(If \(\mathcal{F}\) is irreducible, both
\(\mathrm{H}^0((\mathbb{P}^1)^{\mathrm{an}},j_*\operatorname{End}(\mathcal{F}))\)
and its dual
\(\mathrm{H}^2((\mathbb{P}^1)^{\mathrm{an}},j_*\operatorname{End}(\mathcal{F}))\)
are one-dimensional, so for any nonzero irreducible \(\mathcal{F}\), we have
\(\chi((\mathbb{P}^1)^{\mathrm{an}},j_*\operatorname{End}(\mathcal{F}))\leq 2\).)

We will resort to a transcendental argument. For a suitable choice of base point
\(u\) in \(U^{\mathrm{an}}\), a suitable numbering \(s_1,\ldots,s_k\) of the
\(k:=\operatorname{Card}(S)\) points at \(\infty\), and suitably chosen loops
\(\gamma_i\) which run from \(u\) to \(D^*(s_i)\), turn once counterclockwise
around \(s_i\), and then return to \(u\) the same way they came, the fundamental
group \(\pi_1(U^{\mathrm{an}},u)\) may be described in terms of generators and
relations as the abstract group \(\Gamma_k\) with \(k\) generators \(C_i\)
subject to the one relation
\[
  \prod_i C_i := C_1C_2\cdots C_k=1,
\]
via the isomorphism \(C_i\mapsto\gamma_i\). Notice that the conjugacy class of
\(C_i\) in \(\Gamma_k\) is that of local monodromy \(\gamma(s_i)\) around
\(s_i\).

From this point of view, a rank \(n\) local system \(\mathcal{F}\) on
\(U^{\mathrm{an}}\) is a collection of \(k\) elements \(A_i\) in
\(\operatorname{GL}(n,\mathbb{C})\) which satisfy \(\prod_i A_i=1\). Given a
second rank \(n\) local system \(\mathcal{G}\) on \(U^{\mathrm{an}}\),
corresponding to a collection of \(k\) elements \(D_i\) in
\(\operatorname{GL}(n,\mathbb{C})\) which satisfy \(\prod_i D_i=1\),
\(\mathcal{F}\) and \(\mathcal{G}\) have isomorphic local monodromy if and only
if for each \(i\), \(A_i\) and \(D_i\) are conjugate in
\(\operatorname{GL}(n,\mathbb{C})\), i.e., if and only if for each \(i\) there
exists an element \(B_i\) in \(\operatorname{GL}(n,\mathbb{C})\) such that
\[
  D_i=B_iA_iB_i^{-1}.
\]
\(\mathcal{F}\) and \(\mathcal{G}\) are isomorphic if and only if there exists a
single element \(C\) in \(\operatorname{GL}(n,\mathbb{C})\), or equivalently in
\(\operatorname{SL}(n,\mathbb{C})\), such that
\[
  CA_iC^{-1}=D_i\quad\text{for all }i.
\]

So suppose that \(\mathcal{F}\) is a rank \(n\) local system, corresponding to
a system of \(k\) elements \(A_i\) in \(\operatorname{GL}(n,\mathbb{C})\) which
satisfy \(\prod_i A_i=1\). Then \(\mathcal{F}\) is physically rigid if and only
if given any system of \(k\) elements \(B_i\) in
\(\operatorname{GL}(n,\mathbb{C})\) such that
\[
  \prod_i (B_iA_iB_i^{-1})=1,
\]
there exists a single element \(C\) in \(\operatorname{SL}(n,\mathbb{C})\) such
that
\[
  CA_iC^{-1}=B_iA_iB_i^{-1}\quad\text{for all }i.
\]

Fix such an \(\mathcal{F}\). By the Euler--Poincaré formula, we have
\[
  \chi\bigl((\mathbb{P}^1)^{\mathrm{an}},j_*\operatorname{End}(\mathcal{F})\bigr)
  =
  (2-k)n^2+\sum_i\dim(\mathfrak{z}(A_i)),
\]
where \(\mathfrak{z}(A_i)\) denotes the commuting algebra of \(A_i\) in
\(\operatorname{M}(n,\mathbb{C})\). Let us denote by \(Z(A_i)\) the subgroup of
\(\operatorname{GL}(n,\mathbb{C})\) consisting of all elements which commute
with \(A_i\). Then \(Z(A_i)\) is a nonempty, and hence dense, open set of the
linear space \(\mathfrak{z}(A_i)\), namely it is the open set where the
determinant is invertible. Therefore \(Z(A_i)\) is irreducible, of
\(\dim(Z(A_i))=\dim(\mathfrak{z}(A_i))\). Thus we may rewrite the above formula
as
\[
  \chi\bigl((\mathbb{P}^1)^{\mathrm{an}},j_*\operatorname{End}(\mathcal{F})\bigr)
  =
  (2-k)n^2+\sum_i\dim(Z(A_i)).
\]

Consider the \(k\)-fold self product of \(\operatorname{GL}(n,\mathbb{C})\) with
itself,
\[
  X:=(\operatorname{GL}(n,\mathbb{C}))^k,
\]
and the map
\[
  \pi:X\to\operatorname{SL}(n,\mathbb{C})
\]
defined by
\[
  (B_1,\ldots,B_k)\mapsto\prod_i(B_iA_iB_i^{-1}).
\]
(Since \(\prod_i A_i=1\), this map \(\pi\) does indeed land in
\(\operatorname{SL}(n,\mathbb{C})\).) Consider the group
\[
  G:=\operatorname{SL}(n,\mathbb{C})\times\prod_i Z(A_i),
\]
which acts on \(X\) by having an element \((C,Z_1,\ldots,Z_k)\) in \(G\) act on
\(X\) as
\[
  (B_1,\ldots,B_k)\mapsto(CB_1Z_1^{-1},\ldots,CB_kZ_k^{-1}).
\]
The same group \(G\) acts on \(\operatorname{SL}(n,\mathbb{C})\), by having
\((C,Z_1,\ldots,Z_k)\) in \(G\) act on \(\operatorname{SL}(n,\mathbb{C})\) as
\[
  A\mapsto CAC^{-1}.
\]
With respect to these actions of \(G\), the map
\(\pi:X\to\operatorname{SL}(n,\mathbb{C})\) is easily checked to be
\(G\)-equivariant.

Since the point \(1\) in \(\operatorname{SL}(n,\mathbb{C})\) is a fixed point
of the \(G\)-action, the group \(G\) acts on the fibre \(\pi^{-1}(1)\). The key
tautology is that \(\mathcal{F}\) is physically rigid if and only if the group
\(G\) acts transitively on \(\pi^{-1}(1)\). Suppose now that \(\mathcal{F}\) is
physically rigid. Then \(G\) acts transitively on \(\pi^{-1}(1)\), so we must
have the inequality
\[
  \dim(G)\geq \dim(\pi^{-1}(1)).
\]
Since the point \(1\) in \(\operatorname{SL}(n,\mathbb{C})\) is defined in
\(\operatorname{SL}(n,\mathbb{C})\) by \(n^2-1\) equations,
\(\pi^{-1}(1)\) is defined in \(X\) by \(n^2-1\) equations. Therefore every
irreducible component \(W\) of \(\pi^{-1}(1)\) has
\[
  \dim(W)\geq\dim(X)-(n^2-1).
\]
Since the fibre \(\pi^{-1}(1)\) is nonempty, we have
\[
  \dim(G)\geq\dim(\pi^{-1}(1))
  =
  \sup_{\mathrm{irred\ comp's}\ W}\dim(W)
  \geq
  \dim(X)-(n^2-1)
\]
for \(\mathcal{F}\) physically rigid. Recalling the definitions of \(G\) and
\(X\), the above inequality says
\[
  (n^2-1)+\sum_i\dim(Z(A_i))\geq kn^2-(n^2-1),
\]
i.e.,
\[
  (2-k)n^2+\sum_i\dim(Z(A_i))\geq 2,
\]
i.e.,
\[
  \chi\bigl((\mathbb{P}^1)^{\mathrm{an}},
    j_*\operatorname{End}(\mathcal{F})\bigr)\geq 2.
\]
This completes the proof.
\end{proof}

\setcounter{theorem}{2}
\begin{corollary}
\label{cor:1.1.3}
Notations as in \hyperref[thm:1.1.2]{Theorem~1.1.2} above, let \(\mathcal{F}\) be an irreducible
local system on \(U^{\mathrm{an}}\) of rank \(n\geq 1\), and let
\(\mathcal{L}\) be a rank one local system on \(U^{\mathrm{an}}\). Then the
following conditions are equivalent:
\begin{enumerate}[label=(\arabic*)]
\item \(\mathcal{F}\) is physically rigid.
\item \(\mathcal{F}\otimes\mathcal{L}\) is physically rigid.
\item The dual local system
\(\mathcal{F}^{\vee}:=\operatorname{Hom}(\mathcal{F},\mathbb{C}_{U^{\mathrm{an}}})\)
is physically rigid.
\end{enumerate}
\end{corollary}

\begin{proof}
Indeed, all three local systems \(\mathcal{F}\),
\(\mathcal{F}\otimes\mathcal{L}\) and \(\mathcal{F}^{\vee}\) are irreducible,
and all have the same \(\operatorname{End}\) sheaf on \(U^{\mathrm{an}}\), and
hence the same \(j_*\operatorname{End}\) sheaf on \(\mathbb{P}^1\).
\end{proof}

\section{The case of higher genus}
\label{sec:1.2}

\sourcepara{1.2.1}
Using the same technique, we can analyse the situation in higher genus. We
return to the situation \(X\) a projective smooth connected curve over
\(\mathbb{C}\), of genus \(g\), \(S\) a nonempty finite subset of
\(X(\mathbb{C})\), \(U:=X-S\) the open complement, \(j:U\to X\) the inclusion.
We have already seen that as soon as \(g\geq 1\), no nonzero local system on
\(U^{\mathrm{an}}\) can be physically rigid, due to the possibility of
tensoring with rank one local systems \(\mathcal{L}\) on \(X\). So we introduce
two weaker notions. We say that a local system \(\mathcal{F}\) on
\(U^{\mathrm{an}}\) is weakly physically semi-rigid if there exists a finite
collection of local systems \(\mathcal{F}_1,\mathcal{F}_2,\ldots,\mathcal{F}_d\)
on \(U^{\mathrm{an}}\) with the following property: for any local system
\(\mathcal{G}\) on \(U^{\mathrm{an}}\) such that \(\mathcal{F}\) and
\(\mathcal{G}\) have isomorphic local monodromy, there exists a rank one local
system \(\mathcal{L}\) on \(X\), an index \(1\leq i\leq d\), and an isomorphism
\[
  \mathcal{G}\cong \mathcal{F}_i\otimes j^*\mathcal{L}
\]
of local systems on \(U^{\mathrm{an}}\). We say that a local system
\(\mathcal{F}\) on \(U^{\mathrm{an}}\) is weakly physically rigid if it is
weakly physically semi-rigid and we may take \(d=1\), i.e., if for any local
system \(\mathcal{G}\) on \(U^{\mathrm{an}}\) such that \(\mathcal{F}\) and
\(\mathcal{G}\) have isomorphic local monodromy, there exists a rank one local
system \(\mathcal{L}\) on \(X\) and an isomorphism
\[
  \mathcal{G}\cong\mathcal{F}\otimes j^*\mathcal{L}.
\]

\setcounter{theorem}{1}
\begin{lemma}
\label{lem:1.2.2}
Let \(X\) be a projective smooth connected curve over \(\mathbb{C}\), of genus
\(g\), \(S\) a nonempty finite subset of \(X(\mathbb{C})\), \(U:=X-S\) the open
complement.
\leavevmode\par
\begin{enumerate}[label=(\arabic*)]
\item Any local system \(\mathcal{F}\) on \(U^{\mathrm{an}}\) of rank one is
weakly physically rigid.
\item If \(\mathcal{F}_1\) is a rank one local system on \(U^{\mathrm{an}}\),
then for any nonzero local system \(\mathcal{G}\) on \(U^{\mathrm{an}}\),
\(\mathcal{G}\) is weakly physically rigid if and only if
\(\mathcal{G}\otimes\mathcal{F}_1\) is weakly physically rigid.
\item If \(\mathcal{F}\) is a nonzero local system on \(U^{\mathrm{an}}\),
\(T\) any finite subset of \(U^{\mathrm{an}}(\mathbb{C})\), and
\(k:U^{\mathrm{an}}-T\to U^{\mathrm{an}}\) the inclusion, then
\(\mathcal{F}\) is weakly physically rigid on \(U^{\mathrm{an}}\) if and only
if \(k^*\mathcal{F}\) is weakly physically rigid on \(U^{\mathrm{an}}-T\).
\end{enumerate}
\end{lemma}

\begin{proof}
(1) If \(\mathcal{F}\) and \(\mathcal{G}\) are any two local systems of rank
one with isomorphic local monodromy, then
\(\operatorname{Hom}(\mathcal{G},\mathcal{F})\) is a rank one local system on
\(U^{\mathrm{an}}\) with trivial local monodromy, so of the form \(j^*\mathcal{L}\)
for a unique rank one local system \(\mathcal{L}\), namely
\(j_*\operatorname{Hom}(\mathcal{G},\mathcal{F})\), on \(X^{\mathrm{an}}\),
whence \(\mathcal{F}\cong\mathcal{G}\otimes j^*\mathcal{L}\).

(2) It suffices to show that if \(\mathcal{G}\) is weakly physically rigid then
\(\mathcal{G}\otimes\mathcal{F}_1\) is weakly physically rigid. If
\(\mathcal{H}\) and \(\mathcal{G}\otimes\mathcal{F}_1\) have isomorphic local
monodromy, then \(\mathcal{H}\otimes(\mathcal{F}_1)^{\otimes -1}\) and
\(\mathcal{G}\) have isomorphic local monodromy, so by the weak physical
rigidity of \(\mathcal{G}\), there exists a rank one \(\mathcal{L}\) on
\(X^{\mathrm{an}}\) and an isomorphism
\[
  \mathcal{G}\cong
  (j^*\mathcal{L})\otimes\mathcal{H}\otimes(\mathcal{F}_1)^{\otimes -1}
\]
on \(U^{\mathrm{an}}\). Tensoring with \(\mathcal{F}_1\) gives the required
isomorphism
\[
  \mathcal{G}\otimes\mathcal{F}_1\cong (j^*\mathcal{L})\otimes\mathcal{H}.
\]

(3) Suppose first that \(\mathcal{F}\) is weakly physically rigid on
\(U^{\mathrm{an}}\), and that \(\mathcal{G}\) is a local system on
\(U^{\mathrm{an}}-T\) such that \(k^*\mathcal{F}\) and \(\mathcal{G}\) have
isomorphic local monodromy on \(U^{\mathrm{an}}-T\). Because \(k^*\mathcal{F}\)
has trivial local monodromy at each point of \(T\), so also does
\(\mathcal{G}\), and therefore \(\mathcal{F}(=k_*k^*\mathcal{F})\) and
\(k_*\mathcal{G}\) are two local systems on \(U^{\mathrm{an}}\) with
isomorphic local monodromy. Since \(\mathcal{F}\) is weakly physically rigid on
\(U^{\mathrm{an}}\), there exists a rank one local system \(\mathcal{L}\) on
\(X^{\mathrm{an}}\) and an isomorphism
\[
  \mathcal{F}\cong k_*\mathcal{G}\otimes j^*\mathcal{L}
\]
on \(U^{\mathrm{an}}\). Restricting this isomorphism to \(U^{\mathrm{an}}-T\)
gives
\[
  k^*\mathcal{F}\cong\mathcal{G}\otimes k^*j^*\mathcal{L},
\]
as required.

Conversely, suppose \(\mathcal{F}\) is a nonzero local system on
\(U^{\mathrm{an}}\), such that \(k^*\mathcal{F}\) is weakly physically rigid on
\(U^{\mathrm{an}}-T\). Let \(\mathcal{G}\) be a local system on
\(U^{\mathrm{an}}\) such that \(\mathcal{F}\) and \(\mathcal{G}\) have
isomorphic local monodromy. Then \(k^*\mathcal{F}\) and \(k^*\mathcal{G}\) have
isomorphic local monodromy on \(U^{\mathrm{an}}-T\), so there exists a rank one
local system \(\mathcal{L}\) on \(X^{\mathrm{an}}\) and an isomorphism
\[
  k^*\mathcal{F}\cong k^*\mathcal{G}\otimes k^*j^*\mathcal{L}
  =k^*(\mathcal{G}\otimes j^*\mathcal{L}).
\]
Applying \(k_*\) gives the required isomorphism
\[
  \mathcal{F}=k_*k^*\mathcal{F}
  \cong k_*(k^*(\mathcal{G}\otimes j^*\mathcal{L}))
  =
  \mathcal{G}\otimes j^*\mathcal{L}.
\]
\end{proof}

In a more serious vein, we have:

\setcounter{theorem}{2}
\begin{proposition}
\label{prop:1.2.3}
Let \(X\) be a projective smooth connected curve over \(\mathbb{C}\), of genus
\(g\), \(S\) a nonempty finite subset of \(X(\mathbb{C})\),
\(U:=X-S\) the open complement, \(j:U\to X\) the inclusion. If a nonzero local
system \(\mathcal{F}\) on \(U^{\mathrm{an}}\) is weakly physically semi-rigid,
then
\[
  \chi(X^{\mathrm{an}},j_*\operatorname{End}(\mathcal{F}))\geq 2-2g.
\]
\end{proposition}

\begin{proof}
Once again we resort to a transcendental proof. With suitable base point \(u\)
in \(U^{\mathrm{an}}\), suitable numbering \(s_1,\ldots,s_k\) of the
\(k:=\operatorname{Card}(S)\) points at \(\infty\), and suitable numbering of
the \(g\) ``handles'', the fundamental group \(\pi_1(U^{\mathrm{an}},u)\) may
be described in terms of generators and relations as the abstract group
\(\Gamma_{g,k}\) with \(2g+k\) generators
\[
  E_1,F_1,\ldots,E_g,F_g,C_1,\ldots,C_k
\]
subject to the single relation
\[
  \left(\prod_{j=1,\ldots,g}\{E_j,F_j\}\right)
  \left(\prod_{i=1,\ldots,k} C_i\right)=1,
\]
where we write \(\{a,b\}\) for the commutator \(aba^{-1}b^{-1}\). In this
presentation, the elements \(C_i\) are the local monodromies around the points
\(s_i\) at \(\infty\), and \(\pi_1(X^{\mathrm{an}},u)\) is the quotient of
\(\Gamma_{g,k}\) by the normal subgroup generated by the elements
\(C_1,\ldots,C_k\).

In terms of this presentation of \(\pi_1(U^{\mathrm{an}},u)\), a local system
\(\mathcal{F}\) on \(U^{\mathrm{an}}\) of rank \(n\geq 1\) is a collection
\[
  M_1,N_1,\ldots,M_g,N_g,A_1,\ldots,A_k
\]
which satisfy
\[
  \left(\prod_{j=1,\ldots,g}\{M_j,N_j\}\right)
  \left(\prod_{i=1,\ldots,k} A_i\right)=1.
\]
Fix such an \(\mathcal{F}\). By the Euler--Poincaré formula, we have
\[
  \chi(X^{\mathrm{an}},j_*\operatorname{End}(\mathcal{F}))
  =
  (2-2g-k)n^2+\sum_i\dim(Z(A_i)).
\]
We now repeat the dimension argument of
\hyperref[thm:1.1.2]{Theorem~1.1.2} with
\[
  X=(\operatorname{GL}(n,\mathbb{C}))^{2g+k}
\]
and
\[
  G=\operatorname{SL}(n,\mathbb{C})\times(\mathbb{C}^{\times})^{2g}
    \times\prod_i Z(A_i).
\]
The same equivariance and fibre argument shows that weak physical semi-rigidity
implies
\[
  \dim(G)\geq \dim(\pi^{-1}(1))\geq \dim(X)-(n^2-1).
\]
This says
\[
  (n^2-1)+2g+\sum_i\dim(Z(A_i))\geq (2g+k)n^2-(n^2-1),
\]
i.e.,
\[
  (2-2g-k)n^2+\sum_i\dim(Z(A_i))\geq 2-2g,
\]
i.e.,
\[
  \chi(X^{\mathrm{an}},j_*\operatorname{End}(\mathcal{F}))\geq 2-2g.
\]
\end{proof}

\setcounter{theorem}{3}
\begin{corollary}
\label{cor:1.2.4}
Let \(X\) be a projective smooth connected curve over \(\mathbb{C}\), of genus
\(g\), \(S\) a nonempty finite subset of \(X(\mathbb{C})\), \(k:=\operatorname{Card}(S)\),
\(U:=X-S\) the open complement, \(j:U\to X\) the inclusion. If \(g\geq 2\), and
\(\mathcal{F}\) is a local system on \(U^{\mathrm{an}}\) of rank \(n\geq 2\),
then \(\mathcal{F}\) is not weakly physically semi-rigid.
\end{corollary}

\begin{proof}
For \(\mathcal{F}\) of rank \(n\geq 1\) on \(U^{\mathrm{an}}\), the
Euler--Poincaré formula gives
\begin{align*}
  &\chi(X^{\mathrm{an}},j_*\operatorname{End}(\mathcal{F}))-(2-2g)\\
  &\quad =
  (2-2g-k)n^2+\sum_i\dim(Z(A_i))-(2-2g)\\
  &\quad =
  (2-2g)(n^2-1)+\sum_i(\dim(Z(A_i))-n^2).
\end{align*}
Each term \(\dim(Z(A_i))-n^2\) is \(\leq 0\), with equality if and only if
\(A_i\) is scalar. If \(g\geq 2\) and \(n\geq 2\), then
\((2-2g)(n^2-1)<0\), so
\[
  \chi(X^{\mathrm{an}},j_*\operatorname{End}(\mathcal{F}))<2-2g.
\]
Therefore \(\mathcal{F}\) is not weakly physically semi-rigid.
\end{proof}

\setcounter{theorem}{4}
\begin{corollary}
\label{cor:1.2.5}
(mise pour memoire) Let \(S\) be a nonempty finite subset of
\(\mathbb{P}^1(\mathbb{C})\), \(U:=\mathbb{P}^1-S\) the open complement,
\(j:U\to\mathbb{P}^1\) the inclusion. Let \(\mathcal{F}\) be an irreducible
local system on \(U^{\mathrm{an}}\) of rank \(n\geq 1\). The following
conditions are equivalent:
\begin{enumerate}[label=(\arabic*)]
\item \(\mathcal{F}\) is physically rigid.
\item \(\mathcal{F}\) is weakly physically rigid.
\item \(\mathcal{F}\) is weakly physically semi-rigid.
\item
\(\chi((\mathbb{P}^1)^{\mathrm{an}},j_*\operatorname{End}(\mathcal{F}))\geq 2\).
\item
\(\chi((\mathbb{P}^1)^{\mathrm{an}},j_*\operatorname{End}(\mathcal{F}))=2\).
\end{enumerate}
\end{corollary}

\begin{proof}
The implications \((1)\Rightarrow(2)\Rightarrow(3)\) are trivial, and we have
proven \((3)\Rightarrow(4)\) in
\hyperref[prop:1.2.3]{Proposition~1.2.3} above. For \(\mathcal{F}\) irreducible,
we have already proven \((4)\Rightarrow(5)\) and \((5)\Rightarrow(1)\) in the
proof of \hyperref[thm:1.1.2]{Theorem~1.1.2}.
\end{proof}

\section{The case of genus one}
\label{sec:1.3}

\setcounter{theorem}{0}
\begin{corollary}
\label{cor:1.3.1}
Let \(X\) be a projective smooth connected curve over \(\mathbb{C}\), of genus
\(g=1\), \(S\) a nonempty finite subset of \(X(\mathbb{C})\),
\(U:=X-S\) the open complement, \(j:U\to X\) the inclusion. Let
\(\mathcal{F}\) be a local system on \(U^{\mathrm{an}}\) of rank \(n\geq 1\).
Consider the following conditions:
\leavevmode\par
\begin{enumerate}[label=(\arabic*)]
\item \(\mathcal{F}\) is weakly physically rigid.
\item \(\mathcal{F}\) is weakly physically semi-rigid.
\item \(\chi(X^{\mathrm{an}},j_*\operatorname{End}(\mathcal{F}))\geq 0=(2-2g)\).
\item \(\chi(X^{\mathrm{an}},j_*\operatorname{End}(\mathcal{F}))=0=(2-2g)\).
\item \(\mathcal{F}\) has all its local monodromies scalar.
\item \(j_*\operatorname{End}(\mathcal{F})\) is a local system on \(X^{\mathrm{an}}\).
\end{enumerate}
We have the implications
\[
  (1)\Rightarrow(2)\Rightarrow(3)\Rightarrow(4)\Rightarrow(5)\Rightarrow(6).
\]
If in addition \(\mathcal{F}\) is irreducible, these conditions are all
equivalent.
\end{corollary}

\begin{proof}
The implication \((1)\Rightarrow(2)\) is trivial, and \((2)\Rightarrow(3)\) is
the content of \hyperref[prop:1.2.3]{Proposition~1.2.3} above. If \(g=1\), the Euler--Poincaré formula
gives
\begin{align*}
  &\chi(X^{\mathrm{an}},j_*\operatorname{End}(\mathcal{F}))-(2-2g)\\
  &\quad =
  (2-2g-k)n^2+\sum_i\dim(Z(A_i))-(2-2g)\\
  &\quad =
  \sum_i(\dim(Z(A_i))-n^2).
\end{align*}
Thus
\(\chi(X^{\mathrm{an}},j_*\operatorname{End}(\mathcal{F}))\leq 2-2g=0\),
with strict inequality unless all the local monodromies \(A_i\) are scalar. Thus
\((3)\Rightarrow(4)\) and \((4)\Rightarrow(5)\). If (5) holds, then
\(\operatorname{End}(\mathcal{F})\) has trivial local monodromy, so
\(j_*\operatorname{End}(\mathcal{F})\) is a local system on the elliptic curve
\(X^{\mathrm{an}}\). Thus \((5)\Rightarrow(6)\).

Suppose now that \(\mathcal{F}\) is irreducible, and that (6) holds, i.e.,
\(j_*\operatorname{End}(\mathcal{F})\) is a local system on \(X^{\mathrm{an}}\).
Let \(\mathcal{G}\) be another local system on \(U^{\mathrm{an}}\) such that
\(\mathcal{F}\) and \(\mathcal{G}\) have isomorphic local monodromy. Then
\(\operatorname{Hom}(\mathcal{F},\mathcal{G})\) and
\(\operatorname{End}(\mathcal{F})\) have isomorphic local monodromy. Since
\(j_*\operatorname{End}(\mathcal{F})\) is a local system on \(X^{\mathrm{an}}\),
\(\operatorname{End}(\mathcal{F})\) has trivial local monodromy. Hence
\(\operatorname{Hom}(\mathcal{F},\mathcal{G})\) has trivial local monodromy,
and therefore \(j_*\operatorname{Hom}(\mathcal{F},\mathcal{G})\) is a local
system on \(X^{\mathrm{an}}\). Because the \(\pi_1\) of \(X^{\mathrm{an}}\) is
abelian, any local system on \(X^{\mathrm{an}}\) is a successive extension of
rank one local systems \(\mathcal{L}_k\) on \(X^{\mathrm{an}}\). So there exists
a rank one local system \(\mathcal{L}\) on \(X^{\mathrm{an}}\) and a non-zero
map of local systems on \(X^{\mathrm{an}}\)
\[
  \mathcal{L}\to j_*\operatorname{Hom}(\mathcal{F},\mathcal{G}).
\]
Tensoring this map with \(\mathcal{L}^{\otimes -1}\) gives a nonzero map
\[
  \mathbb{C}\to
  j_*\operatorname{Hom}(\mathcal{F},\mathcal{G})\otimes\mathcal{L}^{\otimes -1}
  =
  j_*\operatorname{Hom}(\mathcal{F}\otimes j^*\mathcal{L},\mathcal{G}),
\]
i.e., a nonzero element in
\[
  \mathrm{H}^0\bigl(X^{\mathrm{an}},
    j_*\operatorname{Hom}(\mathcal{F}\otimes j^*\mathcal{L},\mathcal{G})\bigr)
  =
  \mathrm{H}^0\bigl(U^{\mathrm{an}},
    \operatorname{Hom}(\mathcal{F}\otimes j^*\mathcal{L},\mathcal{G})\bigr)
  =
  \operatorname{Hom}(\mathcal{F}\otimes j^*\mathcal{L},\mathcal{G}).
\]
Because \(\mathcal{F}\), and hence \(\mathcal{F}\otimes j^*\mathcal{L}\), is
irreducible and of the same rank as \(\mathcal{G}\), any such non-zero map of
local systems is an isomorphism. Therefore
\(\mathcal{F}\otimes j^*\mathcal{L}\cong \mathcal{G}\). Thus
\(\mathcal{F}\) is weakly physically rigid, and so \((6)\Rightarrow(1)\) for
\(\mathcal{F}\) irreducible.
\end{proof}

\section{The case of genus one: detailed analysis}
\label{sec:1.4}

\sourcepara{1.4.1}
We now analyze the case of genus one in greater detail. The first step is to
show that up to tensoring with rank one objects, we may reduce to the case when
there is only a single point at \(\infty\).

\setcounter{theorem}{1}
\begin{lemma}
\label{lem:1.4.2}
Let \(X\) be a projective smooth connected curve over \(\mathbb{C}\), of genus
\(g=1\), \(S=\{s_1,s_2,\ldots,s_k\}\) a finite subset of \(X(\mathbb{C})\) with
\(k\geq 2\) points, \(U:=X-S\) the open complement,
\[
  j:U\to X,\qquad j_1:X-S\to X-\{s_1\}
\]
the inclusions. Let \(\mathcal{F}\) be a local system on \(U^{\mathrm{an}}\) of
rank \(n\geq 1\) which is weakly physically rigid. There exists a rank one local
system \(\mathcal{L}\) on \(U^{\mathrm{an}}\), and a weakly physically rigid
local system \(\mathcal{F}_1\) on \((X-\{s_1\})^{\mathrm{an}}\), such that
\[
  \mathcal{F}\cong \mathcal{L}\otimes j_1^*\mathcal{F}_1.
\]
\end{lemma}

\begin{proof}
If \(\mathcal{F}\) on \(U^{\mathrm{an}}\) of rank \(n\geq 1\) is weakly
physically rigid, its local monodromy at each point \(s_i\) in \(S\) is a
scalar, say \(\alpha_i\). Think of \(\pi_1(U^{\mathrm{an}},u)\) as the abstract
group \(\Gamma_{1,k}\) with \(2+k\) generators
\[
  E,F,C_1,\ldots,C_k
\]
subject to the single relation
\[
  \{E,F\}\left(\prod_{i=1,\ldots,k}C_i\right)=1.
\]
The required \(\mathcal{L}\) is any character \(\chi\) of this group which takes
the value \(\alpha_i\) at \(C_i\) for \(i>1\), e.g., one might take
\(\chi(E)=\chi(F)=1\), and
\[
  \chi(C_1):=\left(\prod_{i=2,\ldots,k}\alpha_i\right)^{-1}.
\]
Then \(\mathcal{F}\otimes\mathcal{L}^{\otimes -1}\) has trivial local monodromy
at each of \(s_2,\ldots,s_k\), so it is of the form \(j_1^*\mathcal{F}_1\) for
some local system \(\mathcal{F}_1\) on \((X-\{s_1\})^{\mathrm{an}}\). By
\hyperref[lem:1.2.2]{Lemma~1.2.2(2)}, \(\mathcal{F}\otimes\mathcal{L}^{\otimes -1}\) on \(U^{\mathrm{an}}\)
is weakly physically rigid, and, by \hyperref[lem:1.2.2]{Lemma~1.2.2(3)}, \(\mathcal{F}_1\) is weakly
physically rigid on \((X-\{s_1\})^{\mathrm{an}}\).
\end{proof}

\sourcepara{1.4.3}
We next analyze the irreducible local systems in genus one when there is a
single point at infinity.

\setcounter{theorem}{3}
\begin{lemma}
\label{lem:1.4.4}
Let \(X\) be a projective smooth connected curve over \(\mathbb{C}\), of genus
\(g=1\), \(x_0\) in \(X(\mathbb{C})\) a point, \(U:=X-\{x_0\}\) the open
complement, \(j:U\to X\) the inclusion. Let \(\mathcal{F}\) be a local system
on \(U^{\mathrm{an}}\) of rank \(n\geq 1\). Then the following conditions are
equivalent:
\leavevmode\par
\begin{enumerate}[label=(\arabic*)]
\item \(\mathcal{F}\) is both irreducible and weakly physically rigid.
\item The local monodromy of \(\mathcal{F}\) around \(x_0\) is a scalar
\(\xi\) which is a primitive \(n\)'th root of unity.
\end{enumerate}
\end{lemma}

\begin{proof}
Suppose first that (1) holds. By the weak physical rigidity of
\(\mathcal{F}\), and \hyperref[cor:1.3.1]{Corollary~1.3.1}, \((1)\Rightarrow(5)\), we know that the local
monodromy of \(\mathcal{F}\) around \(x_0\) is a scalar \(\xi\). We first show
that \(\xi\) is necessarily an \(n\)'th root of unity.

Think of \(\pi_1(U^{\mathrm{an}},u)\) as the abstract group \(\Gamma_{1,1}\)
with \(2+1\) generators \(E,F,C\) subject to the single relation
\[
  \{E,F\}C=1,\quad\text{or }C=\{F,E\},
\]
\(C\) being the local monodromy around \(x_0\). Then \(\mathcal{F}\) ``is'' an
\(n\)-dimensional \(\mathbb{C}\)-representation of this group. So its local
monodromy around \(x_0\) lies in \(\operatorname{SL}(n,\mathbb{C})\), because
it is the commutator of two elements of \(\operatorname{GL}(n,\mathbb{C})\).
Being a scalar, \(\xi\) is necessarily an \(n\)'th root of unity.

Given an \(n\)'th root of unity \(\xi\), we construct an explicit
\(n\)-dimensional \(\mathbb{C}\)-representation \(\rho_{n,\xi}\) on
\(V_{n,\xi}\) of \(\pi_1(U^{\mathrm{an}},u)\), which we now think of as the
free group on \(E\) and \(F\), as follows:
\[
  \begin{aligned}
  V_{n,\xi}&=\mathbb{C}[T]/(T^n-1),\\
  \rho_{n,\xi}(E)&\text{ is the automorphism }A:f(T)\mapsto Tf(T),\\
  \rho_{n,\xi}(F)&\text{ is the automorphism }B:f(T)\mapsto f(\xi T).
  \end{aligned}
\]
Clearly \(AB(f)(T)=Tf(\xi T)\), \(BA(f)(T)=\xi Tf(\xi T)\), so
\(BA=\xi AB\), which is to say \(\{B,A\}=\xi\), or equivalently
\(\{A,B\}\xi=1\). Interpret the representation \(\rho_{n,\xi}\) as a local
system \(\mathcal{F}_{n,\xi}\) on \(U^{\mathrm{an}}\). Then its local monodromy
around \(x_0\) is the scalar \(\xi\).

We must show that if \(\mathcal{F}\) is both irreducible and weakly physically
rigid, then \(\xi\) is a primitive \(n\)'th root of unity. We argue by
contradiction. Suppose that for some factorization \(n=dm\), with both \(m\) and
\(d\) integers \(\geq 2\), \(\xi\) were a \(d\)'th root of unity. Then
\(\mathcal{F}\) and \(\bigoplus_m \mathcal{F}_{d,\xi}\) are two local systems on
\(U^{\mathrm{an}}\) with isomorphic local monodromy (namely \(\xi\)). By the
weak physical rigidity of \(\mathcal{F}\), there exists a rank one
\(\mathcal{L}\) on \(X^{\mathrm{an}}\) and an isomorphism
\[
  \mathcal{F}\cong \left(\bigoplus_m\mathcal{F}_{d,\xi}\right)\otimes j^*\mathcal{L}.
\]
But this shows that \(\mathcal{F}\) is in fact reducible. Thus (1) implies (2).

Suppose now that (2) holds. By \hyperref[cor:1.3.1]{Corollary~1.3.1}, it suffices to prove that
\(\mathcal{F}\) is irreducible. Interpreting \(\mathcal{F}\) as a
representation, this results from the following general lemma.
\end{proof}

\setcounter{theorem}{4}
\begin{lemma}
\label{lem:1.4.5}
Let \(K\) be an algebraically closed field, \(n\geq 2\) an integer, \(V\) an
\(n\)-dimensional \(K\)-vector space, \(A\) and \(B\) two elements of
\(\operatorname{GL}(V)\), \(C\) the commutator \(\{B,A\}\). Let
\(\xi_1,\xi_2,\ldots,\xi_n\) be the \(n\) not necessarily distinct eigenvalues
of \(C\). Suppose that for any proper nonempty subset \(S\) of
\(\{1,2,\ldots,n\}\), the product \(\prod_{i\in S}\xi_i\neq 1\). Then the
subgroup of \(\operatorname{GL}(V)\) generated by \(A\) and \(B\) acts
irreducibly on \(V\). In particular, if \(C\) is scalar, equal to a primitive
\(n\)'th root of unity, then the subgroup of \(\operatorname{GL}(V)\) generated
by \(A\) and \(B\) acts irreducibly on \(V\).
\end{lemma}

\begin{proof}
We restate the hypothesis in the following form: for any monic polynomial
\(f(T)\) of degree \(1\leq d<n\) which divides \(\det_V(T-C)\), we have
\[
  \prod_{\text{all roots }\xi\text{ of }f}\xi\neq 1.
\]
We argue by contradiction. If \(W\) is a nontrivial proper subspace of \(V\)
which is mapped to itself by both \(A\) and \(B\), then \(W\) is also mapped to
itself by their inverses, and so \(W\) is mapped to itself by
\(C=\{B,A\}\). But
\[
  C\vert W=\{B\vert W,A\vert W\}
\]
lies in \(\operatorname{SL}(W)\), being a commutator of two elements in
\(\operatorname{GL}(W)\). Taking \(f(T):=\det_W(T-C\vert W)\), we get a
contradiction.
\end{proof}

Using the local systems \(\mathcal{F}_{n,\xi}\) on \((X-\{x_0\})^{\mathrm{an}}\)
constructed in the proof of \hyperref[lem:1.4.4]{Lemma~1.4.4} above, we get a complete description:

\setcounter{theorem}{5}
\begin{proposition}
\label{prop:1.4.6}
Let \(X\) be a projective smooth connected curve over \(\mathbb{C}\), of genus
\(g=1\), \(x_0\) in \(X(\mathbb{C})\) a point, \(U:=X-\{x_0\}\) the open
complement, \(j:U\to X\) the inclusion, \(n\geq 1\) an integer. The local
systems \(\mathcal{F}\) on \(U^{\mathrm{an}}\) of rank \(n\) which are both
irreducible and weakly physically rigid are precisely those of the form
\(\mathcal{F}_{n,\xi}\otimes j^*\mathcal{L}\), with \(\mathcal{L}\) of rank one
on \(X^{\mathrm{an}}\), and with \(\xi\) a primitive \(n\)'th root of unity.
\end{proposition}

\begin{proof}
By \hyperref[lem:1.4.4]{Lemma~1.4.4}, we know if \(\xi\) is a primitive \(n\)'th root of unity, then
\(\mathcal{F}_{n,\xi}\), and hence also
\(\mathcal{F}_{n,\xi}\otimes j^*\mathcal{L}\), is irreducible. Because
\(\mathcal{F}_{n,\xi}\), and hence also
\(\mathcal{F}_{n,\xi}\otimes j^*\mathcal{L}\), have scalar local monodromy,
\hyperref[cor:1.3.1]{Corollary~1.3.1}, \((5)\Rightarrow(1)\), shows that
\(\mathcal{F}_{n,\xi}\otimes j^*\mathcal{L}\) is weakly physically rigid.
Conversely, given a rank \(n\) \(\mathcal{F}\) which is irreducible and weakly
physically rigid, by \hyperref[lem:1.4.4]{Lemma~1.4.4} its local monodromy at \(x_0\) is a scalar \(\xi\)
which is a primitive \(n\)'th root of unity. So \(\mathcal{F}\) and
\(\mathcal{F}_{n,\xi}\) have isomorphic local monodromy. By the weak physical
rigidity of \(\mathcal{F}\), there exists an \(\mathcal{L}\) of rank one on
\(X^{\mathrm{an}}\), and an isomorphism
\[
  \mathcal{F}\cong \mathcal{F}_{n,\xi}\otimes j^*\mathcal{L}.
\]
\end{proof}

\sourcepara{1.4.7}
We next compute the determinant of \(\mathcal{F}_{n,\xi}\). To state the
result, we denote by \(\mathcal{L}_{1/2}\) the rank one local system on
\((X-\{x_0\})^{\mathrm{an}}\) corresponding to the character
\(E\mapsto -1\), \(F\mapsto -1\), and by \(\mathcal{L}_0\) the trivial rank one
local system on \((X-\{x_0\})^{\mathrm{an}}\). (Of course, both of these, like
any rank one local system with only one point at \(\infty\), extend uniquely to
local systems on the complete curve \(X^{\mathrm{an}}\).)

\setcounter{theorem}{7}
\begin{lemma}
\label{lem:1.4.8}
Let \(X\) be a projective smooth connected curve over \(\mathbb{C}\), of genus
\(g=1\), \(x_0\) in \(X(\mathbb{C})\) a point, \(U:=X-\{x_0\}\) the open
complement, \(j:U\to X\) the inclusion, \(n\geq 1\) an integer, \(\xi\) a
primitive \(n\)'th root of unity. Then
\[
  \begin{aligned}
  \det(\mathcal{F}_{n,\xi})&=\mathcal{L}_{1/2}\quad\text{if }n\text{ is even},\\
  \det(\mathcal{F}_{n,\xi})&=\mathcal{L}_0\quad\text{if }n\text{ is odd}.
  \end{aligned}
\]
\end{lemma}

\begin{proof}
As representation, \(\mathcal{F}_{n,\xi}\) is
\[
  \begin{aligned}
  V_{n,\xi}&\text{ is the }n\text{-dimensional }\mathbb{C}\text{-algebra }
  \mathbb{C}[T]/(T^n-1),\\
  \rho_{n,\xi}(E)&\text{ is the automorphism }A:f(T)\mapsto Tf(T),\\
  \rho_{n,\xi}(F)&\text{ is the automorphism }B:f(T)\mapsto f(\xi T).
  \end{aligned}
\]
The assertion is that
\[
  \det(A)=\det(B)=\xi^{n(n+1)/2}\;(=(-1)^{n+1}).
\]
To see this for \(A\), notice that the vectors
\[
  f_i(T):=\sum_{j\bmod n}\xi^{-ij}T^j,\qquad 1\leq i\leq n,
\]
are an eigenbasis for \(A\), with eigenvalues \(\xi^i\). To see it for \(B\),
notice that vectors \(T^i\), \(1\leq i\leq n\), are an eigenbasis for \(B\),
with eigenvalues \(\xi^i\).
\end{proof}

In fact, the local systems \(\mathcal{F}_{n,\xi}\) have a stronger rigidity
property.

\setcounter{theorem}{8}
\begin{proposition}
\label{prop:1.4.9}
Let \(X\) be a projective smooth connected curve over \(\mathbb{C}\), of genus
\(g=1\), \(x_0\) in \(X(\mathbb{C})\) a point, \(U:=X-\{x_0\}\) the open
complement, \(n\geq 1\) an integer. Let \(\mathcal{F}\) be a local system on
\(U^{\mathrm{an}}\) of rank \(n\) which is both irreducible and weakly
physically rigid. Denote by \(\xi\) the primitive \(n\)'th root of unity which
is the local monodromy of \(\mathcal{F}\) around \(x_0\). If
\(\det(\mathcal{F})\cong\det(\mathcal{F}_{n,\xi})\), then
\(\mathcal{F}\cong\mathcal{F}_{n,\xi}\).
\end{proposition}

\begin{proof}
As representation, \(\mathcal{F}_{n,\xi}\) is
\[
  \begin{aligned}
  V_{n,\xi}&\text{ is the }n\text{-dimensional }\mathbb{C}\text{-algebra }
  \mathbb{C}[T]/(T^n-1),\\
  \rho_{n,\xi}(E)&\text{ is the automorphism }A:f(T)\mapsto Tf(T),\\
  \rho_{n,\xi}(F)&\text{ is the automorphism }B:f(T)\mapsto f(\xi T).
  \end{aligned}
\]
By \hyperref[prop:1.4.6]{Proposition~1.4.6}, there exist nonzero scalars
\(\lambda,\mu\) in \(\mathbb{C}^{\times}\)
such that the representation \(\Lambda_{\mathcal{F}}\) corresponding to
\(\mathcal{F}\) is realized on the same space \(V_{n,\xi}\), but with
\[
  \Lambda_{\mathcal{F}}(E):=\lambda A,\qquad
  \Lambda_{\mathcal{F}}(F):=\mu B.
\]
Because \(\det(\mathcal{F})\cong\det(\mathcal{F}_{n,\xi})\), the scalars
\(\lambda\) and \(\mu\) are \(n\)'th roots of unity. Thus we must show that the
automorphisms \(A\) and \(B\) of \(V_{n,\xi}\) have the following property (*):

\par\addvspace{\sourceblockpreskip}%
\noindent(*) for any \(n\)'th roots of unity \(\lambda,\mu\), there exists an
automorphism \(X\) of \(V_{n,\xi}\) such that
\[
  \lambda A=XAX^{-1}\quad\text{and}\quad \mu B=XBX^{-1}.
\]

To prove (*), recall that \(BA=\xi AB\). Thus
\[
  BAB^{-1}=\xi A,\qquad ABA^{-1}=\xi B.
\]
So if we write \(\lambda=\xi^i\) and \(\mu=\xi^j\), we may take
\[
  X=A^jB^i.
\]
\end{proof}

\setcounter{theorem}{9}
\begin{corollary}
\label{cor:1.4.10}
Let \(X\) be a projective smooth connected curve over \(\mathbb{C}\), of genus
\(g=1\), \(x_0\) in \(X(\mathbb{C})\) a point, \(U:=X-\{x_0\}\) the open
complement, \(n\geq 1\) an integer. Let \(\mathcal{F}\) be a local system on
\(U^{\mathrm{an}}\) of rank \(n\) which is both irreducible and weakly
physically rigid. Denote by \(\xi\) the primitive \(n\)'th root of unity which
is the local monodromy of \(\mathcal{F}\) around \(x_0\). The isomorphism class
of \(\mathcal{F}\) is determined by the isomorphism class of the data
\((n,\xi,\det(\mathcal{F}))\).
\end{corollary}

\begin{proof}
The group (under \(\otimes\)) of isomorphism classes of rank one local systems
on \(U^{\mathrm{an}}\) is divisible (being \(\mathbb{C}^{\times}\times
\mathbb{C}^{\times}\)), so tensoring \(\mathcal{F}\) with an \(n\)'th root of
\(\det(\mathcal{F}_{n,\xi})\otimes\det(\mathcal{F})^{-1}\), we reduce to the
proposition above.
\end{proof}

\sourcepara{1.4.11}
We now investigate the situation in genus one when there are \(k\geq 2\) points
at \(\infty\).

\setcounter{theorem}{11}
\begin{proposition}
\label{prop:1.4.12}
Let \(X\) be a projective smooth connected curve over \(\mathbb{C}\), of genus
\(g=1\), \(S=\{s_1,s_2,\ldots,s_k\}\) a finite subset of \(X(\mathbb{C})\) with
\(k\geq 2\) points, \(U:=X-S\) the open complement,
\[
  j:U\to X,\qquad j_1:X-S\to X-\{s_1\}
\]
the inclusions. Fix an integer \(n\geq 1\). The local systems \(\mathcal{F}\) on
\(U^{\mathrm{an}}\) of rank \(n\) which are both irreducible and weakly
physically rigid are precisely those of the form
\[
  (j_1)^*(\mathcal{F}_{n,\xi}\text{ on }(X-\{s_1\})^{\mathrm{an}})\otimes
  \mathcal{L},
\]
with \(\mathcal{L}\) of rank one on \(U^{\mathrm{an}}\), and with \(\xi\) a
primitive \(n\)'th root of unity.
\end{proposition}

\begin{proof}
Simply combine \hyperref[lem:1.4.2]{Lemma~1.4.2} and
\hyperref[prop:1.4.6]{Proposition~1.4.6}.
\end{proof}

Here is an intrinsic characterization.

\setcounter{theorem}{12}
\begin{lemma}
\label{lem:1.4.13}
Let \(X\) be a projective smooth connected curve over \(\mathbb{C}\), of genus
\(g=1\), \(S=\{s_1,\ldots,s_k\}\) a finite subset of \(X(\mathbb{C})\) with
\(k\geq 1\) points, \(U:=X-S\) the open complement, \(j:U\to X\) the inclusion.
Let \(\mathcal{F}\) be a local system on \(U^{\mathrm{an}}\) of rank \(n\geq 1\).
Then \(\mathcal{F}\) is both irreducible and weakly physically rigid if and
only if the following two conditions hold:
\leavevmode\par
\begin{enumerate}[label=(\arabic*)]
\item \(\mathcal{F}\) has scalar monodromy \(\xi_i\) around each point \(s_i\)
in \(S\).
\item \(\prod_{i\in S}\xi_i\) is a primitive \(n\)'th root of unity.
\end{enumerate}
\end{lemma}

\begin{proof}
The assertion is invariant under tensoring with an \(\mathcal{L}\) of rank one
on \(U^{\mathrm{an}}\), so we are reduced, as in the proof of
\hyperref[lem:1.4.2]{Lemma~1.4.2}, to the case when \(k=1\), where it is
\hyperref[lem:1.4.4]{Lemma~1.4.4}.
\end{proof}

\setcounter{theorem}{13}
\begin{proposition}
\label{prop:1.4.14}
Let \(X\) be a projective smooth connected curve over \(\mathbb{C}\), of genus
\(g=1\), \(S=\{s_1,s_2,\ldots,s_k\}\) a finite subset of \(X(\mathbb{C})\) with
\(k\geq 1\) points, \(U:=X-S\) the open complement,
\[
  j:U\to X,\qquad j_1:X-S\to X-\{s_1\}
\]
the inclusions. Let \(\mathcal{F}\) be a local system on \(U^{\mathrm{an}}\) of
rank \(n\geq 1\) which is both irreducible and weakly physically rigid. Denote
by \(\xi_i\) the scalar which is the local monodromy of \(\mathcal{F}\) around
\(s_i\). The isomorphism class of \(\mathcal{F}\) is determined by the
isomorphism class of the data
\[
  (n,\{\xi_i\}_i,\det(\mathcal{F})).
\]
\end{proposition}

\begin{proof}
Given this data, first construct the rank one local system
\(\mathcal{L}(\{\xi_i\}_i)\) on \((X-\{s_1,\ldots,s_k\})^{\mathrm{an}}\) given
as character of \(\pi_1\) by
\[
  E\mapsto 1,\quad F\mapsto 1,\quad
  C_1\mapsto\prod_{i\geq 2}\xi_i,\quad
  C_i\mapsto(\xi_i)^{-1}\quad\text{for }i\geq 2.
\]
Tensoring with \(\mathcal{L}(\{\xi_i\}_i)\), we reduce to the fact
\hyperref[cor:1.4.10]{Corollary~1.4.10} that
\[
  (j_1)_*(\mathcal{F}\otimes\mathcal{L}(\{\xi_i\}_i))
\]
on \((X-\{s_1\})^{\mathrm{an}}\) is determined up to isomorphism by the data
\[
  (n,\prod_i\xi_i,\det((j_1)_*(\mathcal{F}\otimes\mathcal{L}(\{\xi_i\}_i)))).
\]
\end{proof}

\sourcepara{1.4.15}
We now return to the case of a single point at \(\infty\), and give a structure
theorem for irreducible, weakly physically rigid local systems of rank
\(n\geq 1\), which says they are all induced from rank one local systems.

\setcounter{theorem}{15}
\begin{theorem}
\label{thm:1.4.16}
Let \(X\) be a projective smooth connected curve over \(\mathbb{C}\), of genus
\(g=1\), \(x_0\) in \(X(\mathbb{C})\) a point, \(U:=X-\{x_0\}\) the open
complement, \(j:U\to X\) the inclusion, \(n\geq 1\) an integer. Let
\(\mathcal{F}\) be a local system on \(U^{\mathrm{an}}\) of rank \(n\) which is
both irreducible and weakly physically rigid. Denote by \(\xi\) the primitive
\(n\)'th root of unity which is the local monodromy of \(\mathcal{F}\) around
\(x_0\). Let \(\pi:Y\to X\) be any connected finite étale covering of degree
\(n\). Denote by \(S\subset Y(\mathbb{C})\) the set \(\pi^{-1}(x_0)\), which
has \(\operatorname{Card}(S)=n\). Then there exists a rank one local system
\(\mathcal{L}\) on \((Y-S)^{\mathrm{an}}\) for which the local monodromy around
each of the \(n\) points in \(S\) is \(\xi\), and for which
\[
  \pi_*\mathcal{L}\cong\mathcal{F}
\]
on \(U^{\mathrm{an}}\).
\end{theorem}

\begin{proof}
Because \(X\) has genus one, \(\pi_1(X^{\mathrm{an}})\) is abelian,
\(\cong(\mathbb{Z})^2\). So connected finite étale coverings \(Y\) of \(X\)
are in bijective correspondence with the subgroups \(\Gamma\) of
\((\mathbb{Z})^2\) of index \(n\), and such \(\Gamma\) clearly exist. By the
Hurwitz formula, \(Y\) has genus one.

Consider the pullback \(\pi^*\mathcal{F}\) on \((Y-S)^{\mathrm{an}}\). Because
\(\pi\) is finite étale over all of \(X\), the local monodromy of
\(\pi^*\mathcal{F}\) around each point of \(S\) is still the same scalar
\(\xi\). Since \(S\) contains \(n\) points, and \(\xi\) is an \(n\)'th root of
unity, there exist rank one local systems \(\mathcal{L}\) on
\((Y-S)^{\mathrm{an}}\) for which the local monodromy around each of the \(n\)
points in \(S\) is \(\xi\). Pick any such \(\mathcal{L}\). (We will ``correct''
it later.)

The local system \(\operatorname{Hom}(\mathcal{L},\pi^*\mathcal{F})\) on
\((Y-S)^{\mathrm{an}}\) has trivial local monodromy at each point of \(S\). So
denoting by \(k:Y-S\to Y\) the inclusion,
\(k_*\operatorname{Hom}(\mathcal{L},\pi^*\mathcal{F})\) is a local system on
\(Y^{\mathrm{an}}\), which is necessarily a successive extension of rank one
local systems on \(Y^{\mathrm{an}}\). So there exists a rank one
\(\mathcal{L}_0\) on \(Y^{\mathrm{an}}\), and a nonzero element of
\[
  \operatorname{Hom}_{Y^{\mathrm{an}}}
  \bigl(\mathcal{L}_0,k_*\operatorname{Hom}(\mathcal{L},\pi^*\mathcal{F})\bigr)
  =
  \operatorname{Hom}_{(Y-S)^{\mathrm{an}}}
  \bigl(\mathcal{L}\otimes k^*\mathcal{L}_0,\pi^*\mathcal{F}\bigr).
\]
Because \(\mathcal{F}\) is irreducible, and \(\pi\) corresponds to a normal
subgroup of finite index, \(\pi^*\mathcal{F}\) is semisimple. Therefore the
group
\[
  \operatorname{Hom}_{(Y-S)^{\mathrm{an}}}
  \bigl(\pi^*\mathcal{F},\mathcal{L}\otimes k^*\mathcal{L}_0\bigr)
  =
  \operatorname{Hom}_{U^{\mathrm{an}}}
  \bigl(\mathcal{F},\pi_*(\mathcal{L}\otimes k^*\mathcal{L}_0)\bigr)
\]
is nonzero. Since \(\mathcal{F}\) is irreducible, and has the same rank \(n\) as
\(\pi_*(\mathcal{L}\otimes k^*\mathcal{L}_0)\), any nonzero map between them is
an isomorphism. Thus
\[
  \mathcal{F}\cong\pi_*(\mathcal{L}\otimes k^*\mathcal{L}_0).
\]
Since \(\mathcal{L}_0\) was a rank one local system on all of \(Y^{\mathrm{an}}\),
\(\mathcal{L}\) and \(\mathcal{L}\otimes k^*\mathcal{L}_0\) have the same local
monodromy at each point of \(S\). So \(\mathcal{L}\otimes k^*\mathcal{L}_0\)
works as the required ``\(\mathcal{L}\)''.
\end{proof}

\setcounter{theorem}{16}
\begin{remark}
\label{rem:1.4.17}
Here is a variant proof. Consider on \((Y-S)^{\mathrm{an}}\) any
\(\mathcal{L}\) for which the local monodromy around each of the \(n\) points
in \(S\) is \(\xi\). Because \(\pi\) is finite étale over all of \(X\),
\(\pi_*\mathcal{L}\) on \(U^{\mathrm{an}}\) is a rank \(n\) local system which
has its local monodromy around \(x_0\) given by the primitive \(n\)'th root of
unity \(\xi\). By \hyperref[lem:1.4.4]{Lemma~1.4.4}, \(\pi_*\mathcal{L}\) is irreducible and weakly
physically rigid. So by \hyperref[prop:1.4.6]{Proposition~1.4.6}, \(\pi_*\mathcal{L}\) is
\(\mathcal{F}\otimes j^*(\mathcal{L}_1)^{\otimes -1}\) for some local system
\(\mathcal{L}_1\) on \(X^{\mathrm{an}}\), and so
\[
  \mathcal{F}
  \cong
  (\pi_*\mathcal{L})\otimes(j^*\mathcal{L}_1)
  =
  \pi_*(\mathcal{L}\otimes \pi^*j^*\mathcal{L}_1)
  =
  \pi_*(\mathcal{L}\otimes k^*\pi^*\mathcal{L}_1),
\]
and \(\mathcal{L}\otimes k^*\pi^*\mathcal{L}_1\) works as the
``\(\mathcal{L}\)''.
\end{remark}

\sourcepara{1.4.18}
What about reducible local systems which are weakly physically rigid (or even
semi-rigid)?

\setcounter{theorem}{18}
\begin{lemma}
\label{lem:1.4.19}
Let \(X\) be a projective smooth connected curve over \(\mathbb{C}\), of genus
\(g=1\), \(S=\{s_1,s_2,\ldots,s_k\}\) a finite subset of \(X(\mathbb{C})\) with
\(k\geq 1\) points, \(U:=X-S\) the open complement,
\[
  j:U\to X,\qquad j_1:X-S\to X-\{s_1\}
\]
the inclusions. Let \(\mathcal{F}\) be a local system on \(U^{\mathrm{an}}\) of
rank \(n\geq 1\) which is weakly physically semi-rigid. Then \(\mathcal{F}\) is
irreducible (and hence weakly physically rigid, by
\hyperref[cor:1.3.1]{Corollary~1.3.1}).
\end{lemma}

\begin{proof}
We first reduce to the case of a single point at \(\infty\). If \(k\geq 2\),
denote by \(\xi_i\) the scalar which is the local monodromy of \(\mathcal{F}\)
around \(s_i\), and by \(\mathcal{L}(\{\xi_i\}_i)\) the rank one local system
on \((X-\{s_1,\ldots,s_k\})^{\mathrm{an}}\) constructed in the proof of
\hyperref[prop:1.4.14]{Proposition~1.4.14}
above. Then \(\mathcal{F}\otimes\mathcal{L}(\{\xi_i\}_i)\) extends to a local
system on \((X-\{s_1\})^{\mathrm{an}}\) which is still weakly physically
semi-rigid, and it suffices to show that this local system is irreducible on
\((X-\{s_1\})^{\mathrm{an}}\).

We now assume that \(k=1\), \(S=\{s_1\}\), and show that any weakly physically
semi-rigid local system \(\mathcal{F}\) on \((X-\{s_1\})^{\mathrm{an}}\) is
irreducible. Let \(\mathcal{F}\) have rank \(n\geq 1\), and denote by \(\xi\)
its local monodromy around the unique point \(s_1\) at \(\infty\). Then \(\xi\)
is an \(n\)'th root of unity, and we know by
\hyperref[lem:1.4.4]{Lemma~1.4.4} that \(\mathcal{F}\) is
irreducible if and only if \(\xi\) is a primitive \(n\)'th root of unity. So if
\(\xi\) has exact order \(n\), we are done.

So suppose that \(\xi\) has exact order \(d<n\), and define \(m:=n/d\), an
integer \(>1\). Let \(\mathcal{H}\) be any local system on \(X^{\mathrm{an}}\)
of rank \(m\). Then \(\mathcal{F}\) and
\(\mathcal{F}_{d,\xi}\otimes j^*\mathcal{H}\) are two local systems on
\((X-\{s_1\})^{\mathrm{an}}\) with isomorphic local monodromy (namely \(\xi\)).
By the weak physical semi-rigidity of \(\mathcal{F}\), there exists a finite
collection of local systems \(\mathcal{F}_1,\mathcal{F}_2,\ldots,\mathcal{F}_r\)
on \((X-\{s_1\})^{\mathrm{an}}\) with the following property: for any local
system \(\mathcal{G}\) on \(U^{\mathrm{an}}\) such that \(\mathcal{F}\) and
\(\mathcal{G}\) have isomorphic local monodromy, there exists a rank one local
system \(\mathcal{L}\) on \(X\), an index \(1\leq i\leq r\), and an isomorphism
\[
  \mathcal{G}\cong \mathcal{F}_i\otimes j^*\mathcal{L}.
\]
Applying this to
\(\mathcal{G}:=\mathcal{F}_{d,\xi}\otimes j^*\mathcal{H}\), as
\(\mathcal{H}\) runs over all rank \(m\) local systems on \(X^{\mathrm{an}}\),
we see that up to tensoring with a rank one local system on \(X^{\mathrm{an}}\),
we obtain only finitely many isomorphism classes.

We will show this is impossible if \(m>1\). Pick a rank one \(\mathcal{L}_1\) on
\(X^{\mathrm{an}}\) of infinite order. For each integer \(k\geq 1\), define a
rank \(m\) local system \(\mathcal{H}(k)\) on \(X^{\mathrm{an}}\),
\[
  \mathcal{H}(k):=\bigoplus_{1\leq i\leq m}(\mathcal{L}_1)^{\otimes ki}.
\]
Then
\[
  \mathcal{F}_{d,\xi}\otimes j^*\mathcal{H}(k)
  =
  \bigoplus_{1\leq i\leq m}
  \mathcal{F}_{d,\xi}\otimes(j^*\mathcal{L}_1)^{\otimes ki}.
\]
We claim that even modulo tensoring with a rank one local system on
\(X^{\mathrm{an}}\), there are still an infinity of isomorphism classes among
the local systems
\(\mathcal{F}_{d,\xi}\otimes j^*\mathcal{H}(k)\). Notice that each
\(\mathcal{F}_{d,\xi}\otimes j^*\mathcal{H}(k)\) is a direct sum of \(m\)
pairwise non-isomorphic irreducibles each of the same rank \(d\), namely the
\(\mathcal{F}_{d,\xi}\otimes(j^*\mathcal{L}_1)^{\otimes ki}\) with
\(1\leq i\leq m\), whose determinants
\[
  \det(\mathcal{F}_{d,\xi})\otimes(j^*\mathcal{L}_1)^{\otimes dki}
\]
are already pairwise non-isomorphic, \(\mathcal{L}_1\) being of infinite order.

To any direct sum \(\mathcal{G}\) of \(m\) pairwise non-isomorphic irreducibles
\(\mathcal{G}_i\), each of the same rank \(d\), we may attach the following
invariant: the finite set consisting of the distinct isomorphism classes among
the \(m(m-1)\) rank one objects
\[
  \det(\mathcal{G}_i)/\det(\mathcal{G}_j),\quad\text{for all }i\neq j.
\]
This construction visibly attaches the same invariant to \(\mathcal{G}\) and to
\(\mathcal{G}\otimes j^*\mathcal{L}\), for any rank one local system
\(\mathcal{L}\) on \(X^{\mathrm{an}}\).

So it suffices to see that the objects
\[
  \mathcal{F}_{d,\xi}\otimes j^*\mathcal{H}(k)
  =
  \bigoplus_{1\leq i\leq m}
  \mathcal{F}_{d,\xi}\otimes(j^*\mathcal{L}_1)^{\otimes ki},
  \qquad k\geq 1,
\]
each give rise to distinct invariants. But this is obvious. The invariant
attached to \(\mathcal{F}_{d,\xi}\otimes j^*\mathcal{H}(k)\) consists visibly
of the isomorphism classes
\[
  (j^*\mathcal{L}_1)^{\otimes kdp},\qquad -(m-1)\leq p\leq m-1.
\]
Because \(\mathcal{L}_1\) has infinite order,
\((j^*\mathcal{L}_1)^{\otimes kd(m-1)}\), which occurs in the invariant of
\(\mathcal{F}_{d,\xi}\otimes j^*\mathcal{H}(k)\), does not occur in the
invariant of \(\mathcal{F}_{d,\xi}\otimes j^*\mathcal{H}(N)\) for any
\(1\leq N<k\).
\end{proof}
